\documentclass[12pt]{article}
\usepackage[french]{babel}
\usepackage[utf8]{inputenc}
\usepackage[T1]{fontenc}
\usepackage{fancyhdr}
\usepackage{graphicx}

\usepackage{enumerate}
\usepackage{amsmath}
\usepackage{amssymb}
\usepackage{subfigure}
\usepackage{amsmath}
\usepackage{amssymb}
\usepackage{bm}
\usepackage{url}
\usepackage[dvipsnames]{xcolor}
\usepackage{todonotes}


% If your paper is accepted, change the options for the package
% aistats2e as follows:
%
%\usepackage[accepted]{aistats2e}
%
% This option will print headings for the title of your paper and
% headings for the authors names, plus a copyright note at the end of
% the first column of the first page.
\setlength{\parindent}{0cm}
\addtolength{\oddsidemargin}{-2cm}
\addtolength{\evensidemargin}{-2cm}
\setlength{\textwidth}{17.78cm}
\addtolength{\topmargin}{-2.25cm}
\setlength{\textheight}{23.24cm}
\setlength{\headheight}{30pt}
\addtolength{\parskip}{5mm}
\pagestyle{fancy}

%************
%* COMMANDS *
%************

%%%%% NEW MATH DEFINITIONS %%%%%

% Mark sections of captions for referring to divisions of figures
\newcommand{\figleft}{{\em (Left)}}
\newcommand{\figcenter}{{\em (Center)}}
\newcommand{\figright}{{\em (Right)}}
\newcommand{\figtop}{{\em (Top)}}
\newcommand{\figbottom}{{\em (Bottom)}}
\newcommand{\captiona}{{\em (a)}}
\newcommand{\captionb}{{\em (b)}}
\newcommand{\captionc}{{\em (c)}}
\newcommand{\captiond}{{\em (d)}}

% Highlight a newly defined term
\newcommand{\newterm}[1]{{\bf #1}}


% Figure reference, lower-case.
\def\figref#1{figure~\ref{#1}}
% Figure reference, capital. For start of sentence
\def\Figref#1{Figure~\ref{#1}}
\def\twofigref#1#2{figures \ref{#1} and \ref{#2}}
\def\quadfigref#1#2#3#4{figures \ref{#1}, \ref{#2}, \ref{#3} and \ref{#4}}
% Section reference, lower-case.
\def\secref#1{section~\ref{#1}}
% Section reference, capital.
\def\Secref#1{Section~\ref{#1}}
% Reference to two sections.
\def\twosecrefs#1#2{sections \ref{#1} and \ref{#2}}
% Reference to three sections.
\def\secrefs#1#2#3{sections \ref{#1}, \ref{#2} and \ref{#3}}
% Reference to an equation, lower-case.
\def\eqref#1{equation~\ref{#1}}
% Reference to an equation, upper case
\def\Eqref#1{Equation~\ref{#1}}
% A raw reference to an equation---avoid using if possible
\def\plaineqref#1{\ref{#1}}
% Reference to a chapter, lower-case.
\def\chapref#1{chapter~\ref{#1}}
% Reference to an equation, upper case.
\def\Chapref#1{Chapter~\ref{#1}}
% Reference to a range of chapters
\def\rangechapref#1#2{chapters\ref{#1}--\ref{#2}}
% Reference to an algorithm, lower-case.
\def\algref#1{algorithm~\ref{#1}}
% Reference to an algorithm, upper case.
\def\Algref#1{Algorithm~\ref{#1}}
\def\twoalgref#1#2{algorithms \ref{#1} and \ref{#2}}
\def\Twoalgref#1#2{Algorithms \ref{#1} and \ref{#2}}
% Reference to a part, lower case
\def\partref#1{part~\ref{#1}}
% Reference to a part, upper case
\def\Partref#1{Part~\ref{#1}}
\def\twopartref#1#2{parts \ref{#1} and \ref{#2}}

\def\ceil#1{\lceil #1 \rceil}
\def\floor#1{\lfloor #1 \rfloor}
\def\1{\bm{1}}
\newcommand{\train}{\mathcal{D}}
\newcommand{\valid}{\mathcal{D_{\mathrm{valid}}}}
\newcommand{\test}{\mathcal{D_{\mathrm{test}}}}

\def\eps{{\epsilon}}


% Random variables
\def\reta{{\textnormal{$\eta$}}}
\def\ra{{\textnormal{a}}}
\def\rb{{\textnormal{b}}}
\def\rc{{\textnormal{c}}}
\def\rd{{\textnormal{d}}}
\def\re{{\textnormal{e}}}
\def\rf{{\textnormal{f}}}
\def\rg{{\textnormal{g}}}
\def\rh{{\textnormal{h}}}
\def\ri{{\textnormal{i}}}
\def\rj{{\textnormal{j}}}
\def\rk{{\textnormal{k}}}
\def\rl{{\textnormal{l}}}
% rm is already a command, just don't name any random variables m
\def\rn{{\textnormal{n}}}
\def\ro{{\textnormal{o}}}
\def\rp{{\textnormal{p}}}
\def\rq{{\textnormal{q}}}
\def\rr{{\textnormal{r}}}
\def\rs{{\textnormal{s}}}
\def\rt{{\textnormal{t}}}
\def\ru{{\textnormal{u}}}
\def\rv{{\textnormal{v}}}
\def\rw{{\textnormal{w}}}
\def\rx{{\textnormal{x}}}
\def\ry{{\textnormal{y}}}
\def\rz{{\textnormal{z}}}

% Random vectors
\def\rvepsilon{{\mathbf{\epsilon}}}
\def\rvtheta{{\mathbf{\theta}}}
\def\rva{{\mathbf{a}}}
\def\rvb{{\mathbf{b}}}
\def\rvc{{\mathbf{c}}}
\def\rvd{{\mathbf{d}}}
\def\rve{{\mathbf{e}}}
\def\rvf{{\mathbf{f}}}
\def\rvg{{\mathbf{g}}}
\def\rvh{{\mathbf{h}}}
\def\rvu{{\mathbf{i}}}
\def\rvj{{\mathbf{j}}}
\def\rvk{{\mathbf{k}}}
\def\rvl{{\mathbf{l}}}
\def\rvm{{\mathbf{m}}}
\def\rvn{{\mathbf{n}}}
\def\rvo{{\mathbf{o}}}
\def\rvp{{\mathbf{p}}}
\def\rvq{{\mathbf{q}}}
\def\rvr{{\mathbf{r}}}
\def\rvs{{\mathbf{s}}}
\def\rvt{{\mathbf{t}}}
\def\rvu{{\mathbf{u}}}
\def\rvv{{\mathbf{v}}}
\def\rvw{{\mathbf{w}}}
\def\rvx{{\mathbf{x}}}
\def\rvy{{\mathbf{y}}}
\def\rvz{{\mathbf{z}}}

% Elements of random vectors
\def\erva{{\textnormal{a}}}
\def\ervb{{\textnormal{b}}}
\def\ervc{{\textnormal{c}}}
\def\ervd{{\textnormal{d}}}
\def\erve{{\textnormal{e}}}
\def\ervf{{\textnormal{f}}}
\def\ervg{{\textnormal{g}}}
\def\ervh{{\textnormal{h}}}
\def\ervi{{\textnormal{i}}}
\def\ervj{{\textnormal{j}}}
\def\ervk{{\textnormal{k}}}
\def\ervl{{\textnormal{l}}}
\def\ervm{{\textnormal{m}}}
\def\ervn{{\textnormal{n}}}
\def\ervo{{\textnormal{o}}}
\def\ervp{{\textnormal{p}}}
\def\ervq{{\textnormal{q}}}
\def\ervr{{\textnormal{r}}}
\def\ervs{{\textnormal{s}}}
\def\ervt{{\textnormal{t}}}
\def\ervu{{\textnormal{u}}}
\def\ervv{{\textnormal{v}}}
\def\ervw{{\textnormal{w}}}
\def\ervx{{\textnormal{x}}}
\def\ervy{{\textnormal{y}}}
\def\ervz{{\textnormal{z}}}

% Random matrices
\def\rmA{{\mathbf{A}}}
\def\rmB{{\mathbf{B}}}
\def\rmC{{\mathbf{C}}}
\def\rmD{{\mathbf{D}}}
\def\rmE{{\mathbf{E}}}
\def\rmF{{\mathbf{F}}}
\def\rmG{{\mathbf{G}}}
\def\rmH{{\mathbf{H}}}
\def\rmI{{\mathbf{I}}}
\def\rmJ{{\mathbf{J}}}
\def\rmK{{\mathbf{K}}}
\def\rmL{{\mathbf{L}}}
\def\rmM{{\mathbf{M}}}
\def\rmN{{\mathbf{N}}}
\def\rmO{{\mathbf{O}}}
\def\rmP{{\mathbf{P}}}
\def\rmQ{{\mathbf{Q}}}
\def\rmR{{\mathbf{R}}}
\def\rmS{{\mathbf{S}}}
\def\rmT{{\mathbf{T}}}
\def\rmU{{\mathbf{U}}}
\def\rmV{{\mathbf{V}}}
\def\rmW{{\mathbf{W}}}
\def\rmX{{\mathbf{X}}}
\def\rmY{{\mathbf{Y}}}
\def\rmZ{{\mathbf{Z}}}

% Elements of random matrices
\def\ermA{{\textnormal{A}}}
\def\ermB{{\textnormal{B}}}
\def\ermC{{\textnormal{C}}}
\def\ermD{{\textnormal{D}}}
\def\ermE{{\textnormal{E}}}
\def\ermF{{\textnormal{F}}}
\def\ermG{{\textnormal{G}}}
\def\ermH{{\textnormal{H}}}
\def\ermI{{\textnormal{I}}}
\def\ermJ{{\textnormal{J}}}
\def\ermK{{\textnormal{K}}}
\def\ermL{{\textnormal{L}}}
\def\ermM{{\textnormal{M}}}
\def\ermN{{\textnormal{N}}}
\def\ermO{{\textnormal{O}}}
\def\ermP{{\textnormal{P}}}
\def\ermQ{{\textnormal{Q}}}
\def\ermR{{\textnormal{R}}}
\def\ermS{{\textnormal{S}}}
\def\ermT{{\textnormal{T}}}
\def\ermU{{\textnormal{U}}}
\def\ermV{{\textnormal{V}}}
\def\ermW{{\textnormal{W}}}
\def\ermX{{\textnormal{X}}}
\def\ermY{{\textnormal{Y}}}
\def\ermZ{{\textnormal{Z}}}

% Vectors
\def\vzero{{\bm{0}}}
\def\vone{{\bm{1}}}
\def\vmu{{\bm{\mu}}}
\def\vtheta{{\bm{\theta}}}
\def\va{{\bm{a}}}
\def\vb{{\bm{b}}}
\def\vc{{\bm{c}}}
\def\vd{{\bm{d}}}
\def\ve{{\bm{e}}}
\def\vf{{\bm{f}}}
\def\vg{{\bm{g}}}
\def\vh{{\bm{h}}}
\def\vi{{\bm{i}}}
\def\vj{{\bm{j}}}
\def\vk{{\bm{k}}}
\def\vl{{\bm{l}}}
\def\vm{{\bm{m}}}
\def\vn{{\bm{n}}}
\def\vo{{\bm{o}}}
\def\vp{{\bm{p}}}
\def\vq{{\bm{q}}}
\def\vr{{\bm{r}}}
\def\vs{{\bm{s}}}
\def\vt{{\bm{t}}}
\def\vu{{\bm{u}}}
\def\vv{{\bm{v}}}
\def\vw{{\bm{w}}}
\def\vx{{\bm{x}}}
\def\vy{{\bm{y}}}
\def\vz{{\bm{z}}}

% Elements of vectors
\def\evalpha{{\alpha}}
\def\evbeta{{\beta}}
\def\evepsilon{{\epsilon}}
\def\evlambda{{\lambda}}
\def\evomega{{\omega}}
\def\evmu{{\mu}}
\def\evpsi{{\psi}}
\def\evsigma{{\sigma}}
\def\evtheta{{\theta}}
\def\eva{{a}}
\def\evb{{b}}
\def\evc{{c}}
\def\evd{{d}}
\def\eve{{e}}
\def\evf{{f}}
\def\evg{{g}}
\def\evh{{h}}
\def\evi{{i}}
\def\evj{{j}}
\def\evk{{k}}
\def\evl{{l}}
\def\evm{{m}}
\def\evn{{n}}
\def\evo{{o}}
\def\evp{{p}}
\def\evq{{q}}
\def\evr{{r}}
\def\evs{{s}}
\def\evt{{t}}
\def\evu{{u}}
\def\evv{{v}}
\def\evw{{w}}
\def\evx{{x}}
\def\evy{{y}}
\def\evz{{z}}

% Matrix
\def\mA{{\bm{A}}}
\def\mB{{\bm{B}}}
\def\mC{{\bm{C}}}
\def\mD{{\bm{D}}}
\def\mE{{\bm{E}}}
\def\mF{{\bm{F}}}
\def\mG{{\bm{G}}}
\def\mH{{\bm{\text{\textbf{H}}}}}
\def\mI{{\bm{I}}}
\def\mJ{{\bm{J}}}
\def\mK{{\bm{K}}}
\def\mL{{\bm{L}}}
\def\mM{{\bm{M}}}
\def\mN{{\bm{N}}}
\def\mO{{\bm{O}}}
\def\mP{{\bm{P}}}
\def\mQ{{\bm{Q}}}
\def\mR{{\bm{R}}}
\def\mS{{\bm{S}}}
\def\mT{{\bm{T}}}
\def\mU{{\bm{U}}}
\def\mV{{\bm{V}}}
\def\mW{{\bm{W}}}
\def\mX{{\bm{X}}}
\def\mY{{\bm{Y}}}
\def\mZ{{\bm{Z}}}
\def\mBeta{{\bm{\beta}}}
\def\mPhi{{\bm{\Phi}}}
\def\mLambda{{\bm{\Lambda}}}
\def\mSigma{{\bm{\Sigma}}}

% Tensor
\DeclareMathAlphabet{\mathsfit}{\encodingdefault}{\sfdefault}{m}{sl}
\SetMathAlphabet{\mathsfit}{bold}{\encodingdefault}{\sfdefault}{bx}{n}
\newcommand{\tens}[1]{\bm{\mathsfit{#1}}}
\def\tA{{\tens{A}}}
\def\tB{{\tens{B}}}
\def\tC{{\tens{C}}}
\def\tD{{\tens{D}}}
\def\tE{{\tens{E}}}
\def\tF{{\tens{F}}}
\def\tG{{\tens{G}}}
\def\tH{{\tens{H}}}
\def\tI{{\tens{I}}}
\def\tJ{{\tens{J}}}
\def\tK{{\tens{K}}}
\def\tL{{\tens{L}}}
\def\tM{{\tens{M}}}
\def\tN{{\tens{N}}}
\def\tO{{\tens{O}}}
\def\tP{{\tens{P}}}
\def\tQ{{\tens{Q}}}
\def\tR{{\tens{R}}}
\def\tS{{\tens{S}}}
\def\tT{{\tens{T}}}
\def\tU{{\tens{U}}}
\def\tV{{\tens{V}}}
\def\tW{{\tens{W}}}
\def\tX{{\tens{X}}}
\def\tY{{\tens{Y}}}
\def\tZ{{\tens{Z}}}


% Graph
\def\gA{{\mathcal{A}}}
\def\gB{{\mathcal{B}}}
\def\gC{{\mathcal{C}}}
\def\gD{{\mathcal{D}}}
\def\gE{{\mathcal{E}}}
\def\gF{{\mathcal{F}}}
\def\gG{{\mathcal{G}}}
\def\gH{{\mathcal{H}}}
\def\gI{{\mathcal{I}}}
\def\gJ{{\mathcal{J}}}
\def\gK{{\mathcal{K}}}
\def\gL{{\mathcal{L}}}
\def\gM{{\mathcal{M}}}
\def\gN{{\mathcal{N}}}
\def\gO{{\mathcal{O}}}
\def\gP{{\mathcal{P}}}
\def\gQ{{\mathcal{Q}}}
\def\gR{{\mathcal{R}}}
\def\gS{{\mathcal{S}}}
\def\gT{{\mathcal{T}}}
\def\gU{{\mathcal{U}}}
\def\gV{{\mathcal{V}}}
\def\gW{{\mathcal{W}}}
\def\gX{{\mathcal{X}}}
\def\gY{{\mathcal{Y}}}
\def\gZ{{\mathcal{Z}}}

% Sets
\def\sA{{\mathbb{A}}}
\def\sB{{\mathbb{B}}}
\def\sC{{\mathbb{C}}}
\def\sD{{\mathbb{D}}}
% Don't use a set called E, because this would be the same as our symbol
% for expectation.
\def\sF{{\mathbb{F}}}
\def\sG{{\mathbb{G}}}
\def\sH{{\mathbb{H}}}
\def\sI{{\mathbb{I}}}
\def\sJ{{\mathbb{J}}}
\def\sK{{\mathbb{K}}}
\def\sL{{\mathbb{L}}}
\def\sM{{\mathbb{M}}}
\def\sN{{\mathbb{N}}}
\def\sO{{\mathbb{O}}}
\def\sP{{\mathbb{P}}}
\def\sQ{{\mathbb{Q}}}
\def\sR{{\mathbb{R}}}
\def\sS{{\mathbb{S}}}
\def\sT{{\mathbb{T}}}
\def\sU{{\mathbb{U}}}
\def\sV{{\mathbb{V}}}
\def\sW{{\mathbb{W}}}
\def\sX{{\mathbb{X}}}
\def\sY{{\mathbb{Y}}}
\def\sZ{{\mathbb{Z}}}

% Entries of a matrix
\def\emLambda{{\Lambda}}
\def\emA{{A}}
\def\emB{{B}}
\def\emC{{C}}
\def\emD{{D}}
\def\emE{{E}}
\def\emF{{F}}
\def\emG{{G}}
\def\emH{{H}}
\def\emI{{I}}
\def\emJ{{J}}
\def\emK{{K}}
\def\emL{{L}}
\def\emM{{M}}
\def\emN{{N}}
\def\emO{{O}}
\def\emP{{P}}
\def\emQ{{Q}}
\def\emR{{R}}
\def\emS{{S}}
\def\emT{{T}}
\def\emU{{U}}
\def\emV{{V}}
\def\emW{{W}}
\def\emX{{X}}
\def\emY{{Y}}
\def\emZ{{Z}}
\def\emSigma{{\Sigma}}

% entries of a tensor
% Same font as tensor, without \bm wrapper
\newcommand{\etens}[1]{\mathsfit{#1}}
\def\etLambda{{\etens{\Lambda}}}
\def\etA{{\etens{A}}}
\def\etB{{\etens{B}}}
\def\etC{{\etens{C}}}
\def\etD{{\etens{D}}}
\def\etE{{\etens{E}}}
\def\etF{{\etens{F}}}
\def\etG{{\etens{G}}}
\def\etH{{\etens{H}}}
\def\etI{{\etens{I}}}
\def\etJ{{\etens{J}}}
\def\etK{{\etens{K}}}
\def\etL{{\etens{L}}}
\def\etM{{\etens{M}}}
\def\etN{{\etens{N}}}
\def\etO{{\etens{O}}}
\def\etP{{\etens{P}}}
\def\etQ{{\etens{Q}}}
\def\etR{{\etens{R}}}
\def\etS{{\etens{S}}}
\def\etT{{\etens{T}}}
\def\etU{{\etens{U}}}
\def\etV{{\etens{V}}}
\def\etW{{\etens{W}}}
\def\etX{{\etens{X}}}
\def\etY{{\etens{Y}}}
\def\etZ{{\etens{Z}}}

% The true underlying data generating distribution
\newcommand{\pdata}{p_{\rm{data}}}
% The empirical distribution defined by the training set
\newcommand{\ptrain}{\hat{p}_{\rm{data}}}
\newcommand{\Ptrain}{\hat{P}_{\rm{data}}}
% The model distribution
\newcommand{\pmodel}{p_{\rm{model}}}
\newcommand{\Pmodel}{P_{\rm{model}}}
\newcommand{\ptildemodel}{\tilde{p}_{\rm{model}}}
% Stochastic autoencoder distributions
\newcommand{\pencode}{p_{\rm{encoder}}}
\newcommand{\pdecode}{p_{\rm{decoder}}}
\newcommand{\precons}{p_{\rm{reconstruct}}}

\newcommand{\laplace}{\mathrm{Laplace}} % Laplace distribution

\newcommand{\E}{\mathbb{E}}
\newcommand{\Ls}{\mathcal{L}}
\newcommand{\R}{\mathbb{R}}
\newcommand{\emp}{\tilde{p}}
\newcommand{\lr}{\alpha}
\newcommand{\reg}{\lambda}
\newcommand{\rect}{\mathrm{rectifier}}
\newcommand{\softmax}{\mathrm{softmax}}
\newcommand{\sigmoid}{\sigma}
\newcommand{\softplus}{\zeta}
\newcommand{\KL}{D_{\mathrm{KL}}}
\newcommand{\Var}{\mathrm{Var}}
\newcommand{\standarderror}{\mathrm{SE}}
\newcommand{\Cov}{\mathrm{Cov}}
% Wolfram Mathworld says $L^2$ is for function spaces and $\ell^2$ is for vectors
% But then they seem to use $L^2$ for vectors throughout the site, and so does
% wikipedia.
\newcommand{\normlzero}{L^0}
\newcommand{\normlone}{L^1}
\newcommand{\normltwo}{L^2}
\newcommand{\normlp}{L^p}
\newcommand{\normmax}{L^\infty}

\newcommand{\parents}{Pa} % See usage in notation.tex. Chosen to match Daphne's book.

\DeclareMathOperator*{\argmax}{arg\,max}
\DeclareMathOperator*{\argmin}{arg\,min}

\DeclareMathOperator{\sign}{sign}
\DeclareMathOperator{\Tr}{Tr}
\let\ab\allowbreak


%%% new
\newcommand{\diag}{\mathop{\mathrm{diag}}\nolimits}
\usepackage{cancel}
\newcommand{\ttt}[1]{\texttt{#1}}

\newif\ifexercise
\exercisetrue
%\exercisefalse

\newif\ifsolution
\solutiontrue
%\solutionfalse

% Solution blocks (blue), like in the English version
\usepackage{amsthm}
\theoremstyle{definition}
\newtheorem{answer}{Answer}
\definecolor{answer}{rgb}{0.00, 0.12, 0.60}
\makeatletter
\let\answer@orig\answer
\let\endanswer@orig\endanswer
\renewenvironment{answer}{\color{answer}\answer@orig}{\endanswer@orig}
\makeatother
\newcommand{\Answer}[1]{%
  \ifsolution
    {\color{answer}#1}%
  \fi
}

\usepackage{booktabs} % for \toprule, \midrule, \bottomrule
\usepackage{tabularx}
\usepackage[ruled,vlined]{algorithm2e}
\newcommand{\blank}{\underline{\hspace{2.1cm}}}

\usepackage{tikz}
\usetikzlibrary{positioning}

\usepackage{enumitem}
\newcommand{\staritem}{
\addtocounter{enumi}{1}
\item[$\phantom{x}^{*}$\theenumi]}
\setlist[enumerate,1]{leftmargin=*, label=\arabic*.}

% Question/answer headers (consistent with the English version)
\newcommand{\question}[2]{\par\medskip\noindent\textbf{Question #1. #2}\par\medskip}

% Special tokens
\newcommand{\bos}{\texttt{<BOS>}}

\begin{document}


\fancyhead{}
\fancyfoot{}

\fancyhead[L]{
  \begin{tabular}[b]{l}
    IFT6135-A2026  \\
    Prof: Aishwarya Agrawal \\
  \end{tabular}
}
\fancyhead[R]{
  \begin{tabular}[b]{r}
    Devoir 2 - Théorie  \\
    Réseaux de neurones récurrents et Transformers \\
  \end{tabular}
}
\fancyfoot[C]{- Ne pas distribuer -}

\vspace{1cm}

\shorthandoff{:}
{\textbf{Due Date: 15 Novembre 2026 à 23:00 }}\\


\vspace{-0.5cm}
\underline{Instructions}%
\renewcommand{\labelitemi}{\textbullet}

\renewcommand{\labelitemi}{\textbullet}

\begin{itemize}

\item \emph{Ce devoir demande un travail conséquent. Commencez-le suffisamment à l'avance !}
\item \emph{Pour chaque question, détaillez votre raisonnement.}
\item \emph{Rédigez vos réponses en LaTeX en utilisant le modèle fourni. Vous pouvez réutiliser la plupart des notations abrégées, équations et tableaux. Consultez les règles relatives aux devoirs sur le site du cours pour plus de détails.}
\item \emph{Remettez ce devoir sur Gradescope. Veillez à sélectionner les pages contenant la réponse à chaque question. Les questions pour lesquelles aucune page n'a été sélectionnée ne seront pas corrigées.}
\item \emph{Les auxiliaires d'enseignement responsables de ce devoir sont \textbf{Achille Sowa} et \textbf{Djalil Chikhi}.}
\end{itemize}
\vspace{0.2cm}


\vspace{0.2cm}




%\documentclass[10pt,letterpaper]{article}
%\usepackage[utf8]{inputenc}
%\usepackage[T1]{fontenc}
%\usepackage[french]{babel}
%\usepackage{lmodern,amsmath,amssymb,mathtools,booktabs,array,tikz}
%\usepackage[left=2.2cm,right=2.1cm,top=2.55cm,bottom=1.9cm,headheight=29pt,headsep=13pt,footskip=25pt]{geometry}
%\usepackage{fancyhdr,enumitem,xcolor,hyperref}

%\hypersetup{colorlinks=true,urlcolor=blue}
%\setlength{\parindent}{0pt}
%\setlength{\parskip}{8pt}
%\setlist[enumerate]{leftmargin=20pt,itemsep=7pt,topsep=5pt}
%\pagestyle{fancy}\fancyhf{}
%\fancyhead[L]{\begin{tabular}{@{}l@{}}IFT6135-A2026\\Prof. : Aishwarya Agrawal\end{tabular}}
%\fancyhead[R]{\begin{tabular}{@{}r@{}}Devoir 2 -- Théorie\\Réseaux de neurones récurrents et Transformers\end{tabular}}
%\renewcommand{\headrulewidth}{0.4pt}
%\fancyfoot[C]{-- Ne pas distribuer --}
%\renewcommand{\footrulewidth}{0pt}
%\newcommand{\question}[2]{\par\medskip\noindent\textbf{Question #1. #2}\par\medskip}
%\newcommand{\answer}[1]{\par\medskip\noindent\textcolor{blue!65!black}{\textbf{Réponse #1.}}\par\medskip}
%\newcommand{\bos}{\texttt{<BOS>}}
%\newcommand{\R}{\mathbb{R}}
%\begin{document}
%\textbf{Date limite :} le 15 novembre 2026 à 23 h 00

%\textbf{Consignes}
%\begin{itemize}[leftmargin=17pt,itemsep=2pt,topsep=2pt]
%\item Ce devoir demande un travail conséquent. Commencez-le suffisamment à l'avance !
%\item Pour chaque question, détaillez votre raisonnement.
%\item Rédigez vos réponses en LaTeX en utilisant le modèle fourni. Vous pouvez réutiliser la plupart des notations abrégées, équations et tableaux. Consultez les règles relatives aux devoirs sur le site du cours pour plus de détails.
%\item Remettez ce devoir sur Gradescope. Veillez à sélectionner les pages contenant la réponse à chaque question. Les questions pour lesquelles aucune page n'a été sélectionnée ne seront pas corrigées.
%\item Les auxiliaires d'enseignement responsables de ce devoir sont \textbf{Achille Sowa} et \textbf{Djalil Chikhi}.
%\end{itemize}

\question{1 (12)}{Lien entre Transformer linéaire et RNN}\label{ex:lstm}
Soit $\vx_1,\ldots,\vx_L\in\R^d$ une séquence de $L$ vecteurs d'entrée. Un Transformer décodeur utilise une attention softmax causale : au pas de temps $t$, la sortie $\vo_t$ porte son attention sur toutes les positions $i\le t$ de la séquence,
\begin{equation}
\vq_t=\mW_Q\vx_t,\qquad \vk_t=\mW_K\vx_t,\qquad \vv_t=\mW_V\vx_t,\qquad
 \vo_t=\sum_{i=1}^{t}\frac{\exp(\vk_i^\top \vq_t)}{\sum_{j=1}^{t}\exp(\vk_j^\top \vq_t)}\vv_i.
\tag{1}
\end{equation}
 où $\mW_Q,\mW_K\in\R^{d\times d}$ et $\mW_V\in\R^{d\times d}$ sont des matrices de projection apprises. Ainsi, $\vq_t,\vk_t\in\R^d$, $\vv_t\in\R^d$ et $\vo_t\in\R^d$.

L'attention linéaire remplace la similarité $\exp(\vk_i^\top \vq_t)$ par un produit scalaire dans un espace de caractéristiques, $\phi(\vk_i)^\top\phi(\vq_t)$, pour une application de caractéristiques $\phi:\R^{d_k}\to\R^{d_\phi}$. Dans cet exercice, nous considérons le cas le plus simple, $\phi=\mathrm{id}$, c'est-à-dire que la similarité est simplement $\vk_i^\top \vq_t$.

\begin{enumerate}
\item \textbf{(2 points)} Montrez que, lorsque $\exp(\vk_i^\top \vq_t)$ est remplacé par $\vk_i^\top \vq_t$ dans l'équation (1), la sortie peut s'écrire
\begin{equation}
 \vo_t=\frac{\mS_t\,\vq_t}{\vz_t^\top \vq_t},\qquad\text{avec}\qquad \mS_t=\sum_{i=1}^t \vv_i \vk_i^\top\in\R^{d\times d},\qquad \vz_t=\sum_{i=1}^t \vk_i\in\R^d.
\tag{2}
\end{equation}
Cette expression peut être interprétée comme un réseau de neurones récurrent (RNN) dont l'état caché $\mS_t$ est mis à jour à chaque pas de temps, et à partir duquel la sortie est obtenue par projection linéaire avec $\vq_t$. En pratique, le dénominateur est souvent omis, de sorte que la sortie devient
\begin{equation}\vo_t=\mS_t\,\vq_t.\tag{3}\end{equation}
C'est la convention adoptée pour la suite de l'exercice.
\end{enumerate}
%\newpage
\begin{enumerate}[start=2]
\item \textbf{(3 points) Forme récurrente.} À partir de l'équation (3), écrivez un algorithme qui calcule toutes les sorties $o_1,\ldots,o_L$ à partir des entrées $x_1,\ldots,x_L$ en un temps linéaire en la longueur $L$ de la séquence.

\textit{Indication :} établissez d'abord une relation de récurrence exprimant $S_t$ en fonction de $S_{t-1}$, puis utilisez-la pour écrire votre algorithme. Complétez les parties manquantes du squelette ci-dessous.

\begin{minipage}{\linewidth}
    \begin{algorithm}[H]
    \DontPrintSemicolon
    \KwIn{vecteurs $\vx_1, \dots, \vx_L \in \R^{d}$, poids $\mW_Q, \mW_K \in \R^{d \times d}$ et $\mW_V \in \R^{d \times d}$}
    \KwOut{sorties $\vo_1, \dots, \vo_L \in \R^{d}$}
    Initialisation : $\mS_{0} \leftarrow \blank$\;
    \For{$t = 1$ \KwTo $L$}{
      Calculer la requête, la clé et la valeur : $\vq_t \leftarrow \mW_Q \vx_t$, $\vk_t \leftarrow \mW_K \vx_t$, $\vv_t \leftarrow \mW_V \vx_t$\;
      \tcp*[l]{Complétez cette partie}
    }
    \caption{Attention linéaire, forme récurrente}
    \label{alg:linear-rnn}
    \end{algorithm}
\end{minipage}



\item \textbf{(4 points) Forme parallèle.} Sur un GPU, il est plus efficace d'effectuer simultanément de nombreux produits vectoriels sous la forme d'un seul produit matriciel. Empilez les entrées et les sorties comme les lignes de matrices,
\[
\mX=\begin{bmatrix}\vx_1^\top\\\vdots\\\vx_L^\top\end{bmatrix}\in\R^{L\times d},\qquad
\mO=\begin{bmatrix}\vo_1^\top\\\vdots\\\vo_L^\top\end{bmatrix}\in\R^{L\times d},
\]
et posez $\mQ=\mX\mW_Q^\top\in\R^{L\times d}$, $\mK=\mX\mW_K^\top\in\R^{L\times d}$ et $\mV=\mX\mW_V^\top\in\R^{L\times d}$. Montrez que les sorties de l'équation (3) peuvent s'écrire sous forme matricielle comme suit :
\begin{equation}\mO=\bigl(\mQ\mK^\top\odot \mM\bigr)\mV,\tag{4}\end{equation}
où $\odot$ désigne le produit terme à terme (produit de Hadamard) et $\mM\in\{0,1\}^{L\times L}$ est le masque causal, avec $\mM_{ti}=1$ si $i\le t$ et $\mM_{ti}=0$ sinon.

\item \textbf{(3 points) Complexité.} Donnez les complexités en temps et le nombre moyen d'étapes séquentielles pour le calcul de toutes les sorties $o_1,\ldots,o_L$ avec chacune des trois méthodes ci-dessous, en fonction de $L$ et $d$ (vous pouvez ignorer le coût du calcul de $Q$, $K$ et $V$). Justifiez brièvement vos réponses et commentez leur dépendance en $L$.
{\small
\begin{center}
\begin{tabularx}{\linewidth}{@{}Xcc@{}}\toprule
\textbf{Méthode} & \textbf{Complexité en temps} & \textbf{Opérations séquentielles}\\\midrule
Attention softmax (Transformer classique) & & \\
RNN linéaire, forme récurrente (question 2) & & \\
RNN linéaire, forme parallèle, équation (4) (question 3) & & \\\bottomrule
\end{tabularx}
\end{center}
}
\end{enumerate}

\Answer{

}
% Q2


\question{2 (1-1-4-6-2)}{Décodage}
Supposons que nous disposions d'un vocabulaire contenant $N$ mots possibles, dont un jeton spécial $\bos$ indiquant le début d'une phrase. Rappelons qu'en général, un modèle de langage à contexte complet peut s'écrire
\[
p(w_1,w_2,\ldots,w_T\mid w_0)=\prod_{t=1}^T p(w_t\mid w_0,\ldots,w_{t-1}).
\]
%\newpage
Nous utilisons la notation $w_{0:t-1}$ pour désigner la séquence (partielle) $(w_0,\ldots,w_{t-1})$. Une fois notre modèle de langage entièrement entraîné, nous souhaitons générer des séquences de mots réalistes à partir de ce modèle, en commençant par le jeton spécial $\bos$. Nous pouvons notamment chercher à générer la séquence $\vw_{1:T}^*$ la plus probable selon le modèle, définie par
\begin{equation}\vw_{1:T}^*=\arg\max_{\vw_{1:T}}p(\vw_{1:T}\mid w_0=\bos).\tag{5}\end{equation}
Par souci de clarté, nous omettrons désormais le conditionnement explicite sur $w_0$, en supposant que les séquences commencent toujours par le jeton $\bos$.

\textbf{2.1} Combien de séquences possibles de longueur $T+1$ commençant par le jeton $\bos$ peut-on générer au total ? Donnez une expression exacte, sans utiliser la notation $O$. Notez que la longueur $T+1$ inclut ici le jeton $\bos$.

\textbf{2.2} Pour déterminer la séquence $\vw_{1:T}^*$ la plus probable, on pourrait énumérer exhaustivement toutes les séquences possibles et retenir celle dont la probabilité conjointe est la plus élevée (définie dans l'équation (5)). Commentez la faisabilité de cette approche. Passe-t-elle à l'échelle lorsque la taille du vocabulaire augmente ?

\textbf{2.3} Face aux difficultés de l'énumération exhaustive, il est essentiel de recourir à des stratégies de recherche pratiques afin d'obtenir une approximation raisonnable de la séquence la plus probable. Pour générer $B$ séquences de forte vraisemblance, on peut utiliser un algorithme heuristique appelé \textit{recherche en faisceau} (\textit{beam search}), dont le pseudo-code est donné dans l'algorithme 2 ci-dessous.

\begin{minipage}{\linewidth}
    \begin{algorithm}[H]
    \DontPrintSemicolon
    \KwIn{un modèle de langage $p(\vw_{1:T} \mid w_{0})$, la largeur de faisceau (\textit{beam width}) $B$}
    \KwOut{$B$ séquences $\vw_{1:T}^{(b)}$ pour $b \in \{1, \ldots, B\}$}
    Initialisation : $w_{0}^{(b)} \leftarrow \bos$ pour tout $b \in \{1,\ldots,B\}$\;
    Log-vraisemblances initiales : $l_{0}^{(b)} \leftarrow 0$ pour tout $b \in \{1,\ldots,B\}$\;
    \For{$t = 1$ \KwTo $T$}{
      \For{$b = 1$ \KwTo $B$}{
        \For{$j = 1$ \KwTo $N$}{
          $s_{b}(j) \leftarrow l_{t-1}^{(b)} + \log p(w_{t} = j \mid \vw_{0:t-1}^{(b)})$\;
        }
      }
      \For{$b = 1$ \KwTo $B$}{
        Trouver $(b', j)$ tel que $s_{b'}(j)$ soit le $b$-ième score le plus élevé\;
        Sauvegarder la séquence partielle $b'$ : $\widetilde{\vw}_{0:t-1}^{(b)} \leftarrow \vw_{0:t-1}^{(b')}$\;
        Ajouter le mot $j$ à la séquence $b$ : $w_{t}^{(b)} \leftarrow j$\;
        Mettre à jour la log-vraisemblance : $l_{t}^{(b)} \leftarrow s_{b'}(j)$\;
      }
      Affecter les séquences partielles : $\vw_{0:t-1}^{(b)} \leftarrow \widetilde{\vw}_{0:t-1}^{(b)}$ pour tout $b \in \{1, \ldots, B\}$\;
    }
    \caption{Décodage par \textit{Beam search}}
    \label{alg:beam-search-decoding}
    \end{algorithm}
\end{minipage}

    

Quelles sont les complexités en temps et en mémoire de l'algorithme 2 ? Exprimez-les en notation $O$, en fonction de $T$, $B$ et $N$. Cet algorithme de décodage est-il utilisable en pratique lorsque le vocabulaire est très grand ?

\textbf{2.4} Les différentes séquences qu'un modèle de langage peut générer peuvent être représentées par un arbre, dont les nœuds correspondent à des mots et les arêtes sont étiquetées par la log-probabilité
%\newpage
$\log p(w_t\mid \vw_{0:t-1})$, qui dépend du chemin $\vw_{0:t-1}$. Pour cette question, considérez le modèle de langage suivant (les chemins de faible probabilité ont été supprimés pour plus de clarté) :

\hspace*{-2.5em}\begin{minipage}{\linewidth}
\centering
\begin{tikzpicture}[scale=0.9,every node/.style={shape=rectangle, inner sep=5pt, align=center}, level 3/.style={sibling distance=5.5em}, level 2/.style={sibling distance=11em}, level 1/.style={sibling distance=22em}, edge from parent path={(\tikzparentnode.south) .. controls +(0,-0.7) and +(0,0.7) .. (\tikzchildnode.north)}, edge from parent/.style={draw, thick, -latex}, logl/.style={draw=none, font=\footnotesize}, label/.style={anchor=west, align=flush left}]
  \node {\bos}
    child { node {\texttt{the}}
      child { node {\texttt{green\vphantom{l}}}
        child { node {\texttt{ellipse\vphantom{q}}} edge from parent node[above left,logl] {$-3.6$} }
        child { node {\texttt{rectangle}} edge from parent node[above right,logl] {$-3.3$} } edge from parent node[above left,logl] {$-3.4$} }
      child { node {\texttt{red}}
        child { node {\texttt{circle\vphantom{q}}} edge from parent node[above left,logl] {$-4.8$} }
        child { node {\texttt{rectangle}} edge from parent node[above right,logl] {$-4.5$} } edge from parent node[above right,logl] {$-3.1$} } edge from parent node[above left,logl] {$-1.5$} }
    child { node {\texttt{a}}
      child { node {\texttt{blue}}
        child { node {\texttt{ellipse\vphantom{q}}} edge from parent node[above left,logl] {$-5.4$} }
        child { node {\texttt{square\vphantom{l}}} edge from parent node[above right,logl] {$-5.0$} } edge from parent node[above left,logl] {$-2.5$} }
      child { node {\texttt{red}}
        child { node {\texttt{circle\vphantom{q}}} edge from parent node[above left,logl] {$-4.9$} }
        child { node {\texttt{square\vphantom{l}}} edge from parent node[above right,logl] (logl_lvl3) {$-5.1$} } edge from parent node[above right,logl] (logl_lvl2) {$-2.7$} } edge from parent node[above right,logl] (logl_lvl1) {$-1.8$} };

  \node[label, right=1em of logl_lvl3] (logl_lvl3_label) {$\log p(w_{3}\mid \vw_{0:2})$};
  \node[label] (logl_lvl2_label) at (logl_lvl2 -| logl_lvl3_label.west) {$\log p(w_{2}\mid \vw_{0:1})$};
  \node[label] (logl_lvl1_label) at (logl_lvl1 -| logl_lvl3_label.west) {$\log p(w_{1}\mid w_{0})$};
\end{tikzpicture}
\end{minipage}

\textbf{2.4.a} Le \textit{décodage glouton} (\textit{greedy decoding}) est un algorithme simple qui choisit le mot suivant $\overline{w}_t$ en maximisant la probabilité conditionnelle $p(w_t\mid \overline{\vw}_{0:t-1})$ (avec $w_0=\bos$) :
\[
\overline{w_t}=\arg\max_{w_t}\log p(w_t\mid \overline{\vw}_{0:t-1}).
\]
Déterminez $\overline{\vw}_{1:3}$ par décodage glouton pour ce modèle de langage, ainsi que sa log-vraisemblance $\log p(\overline{\vw}_{1:3})$.

\textbf{2.4.b} Appliquez le decodage \textit{Beam search}  avec une largeur $B=2$ à ce modèle de langage et déterminez $\vw_{1:3}^{(1)}$ et $\vw_{1:3}^{(2)}$, ainsi que leurs log-vraisemblances respectives.

\textbf{2.5} Pour chacune des méthodes étudiées (décodage glouton et recherche en faisceau), indiquez la principale limite qui vous semble ressortir.

\Answer{

}
% Q3

\question{3 (10-8-12)}{Température dans l'attention souple}
Considérez le mécanisme d'attention souple (\textit{soft attention}) étudié en cours. Soient $h_1,\ldots,h_{T_x}$ les états de l'encodeur pour une séquence d'entrée de longueur $T_x$. À l'étape de décodage $i$, le décodeur attribue un score d'alignement $e_{ij}$ à chaque état $h_j$ de l'encodeur. Le poids d'attention $\alpha_{ij}$ correspondant détermine la contribution de $h_j$ au vecteur de contexte $c_i$.

Nous modifions les poids d'attention en introduisant une température positive $\tau_i>0$ :
\[
\alpha_{ij}(\tau_i)=\frac{\exp(e_{ij}/\tau_i)}{\sum_{k=1}^{T_x}\exp(e_{ik}/\tau_i)},\qquad c_i=\sum_{j=1}^{T_x}\alpha_{ij}(\tau_i)h_j.
\]
Ici, $i$ désigne l'étape de décodage, $j$ une position dans la séquence d'entrée et $T_x$ la longueur de celle-ci. Les poids d'attention vérifient $\sum_{j=1}^{T_x}\alpha_{ij}=1$.

Dans les questions 1 et 2, nous considérons une étape de décodage fixe. Pour simplifier les notations, posons
\[
e_j:=e_{ij},\qquad\alpha_j:=\alpha_{ij},\qquad\tau:=\tau_i,\qquad c:=c_i.
\]
%\newpage
\begin{enumerate}
\item \textbf{Effet de la température.}
\begin{enumerate}[label=(\alph*),leftmargin=16pt]
\item Considérez les scores d'alignement $e=[2,1,0]$. Calculez $\alpha(1)$ et déterminez
\[\lim_{\tau\to0^+}\alpha(\tau)\qquad\text{et}\qquad\lim_{\tau\to\infty}\alpha(\tau).\]
Donnez les vecteurs de contexte limites correspondants et expliquez brièvement l'effet de $\tau$ sur la sélectivité de l'attention.
\item Pour des scores d'alignement quelconques, montrez que
\[\frac{\partial\alpha_j}{\partial e_k}=\frac{1}{\tau}\alpha_j(\delta_{jk}-\alpha_k),\]
où $\delta_{jk}=1$ si $j=k$, et $0$ sinon.
\item En utilisant le résultat précédent, montrez que
\[\frac{\partial c}{\partial e_k}=\frac{\alpha_k}{\tau}(h_k-c).\]
Expliquez ce que cette expression nous apprend sur l'effet d'une augmentation de $e_k$ sur le vecteur de contexte.
\end{enumerate}
\item \textbf{Changement d'échelle des scores d'alignement.}
\begin{enumerate}[label=(\alph*),leftmargin=16pt]
\item Supposons que tous les scores d'alignement soient multipliés par une même constante $s>0$, de sorte que $e'_j=se_j$, tandis que la température passe de $\tau$ à $\tau'$. Trouvez une relation entre $s$, $\tau$ et $\tau'$ telle que
\[\alpha'_j(\tau')=\alpha_j(\tau)\]
pour tout $j$.
\item Si l'échelle des scores d'alignement et la température sont toutes deux apprises, le fait d'observer $\tau=0.2$ suffit-il pour conclure que le modèle utilise une attention très sélective ? Justifiez votre réponse.
\end{enumerate}
\item \textbf{Température adaptative.} Revenons aux notations initiales, dans lesquelles $i$ désigne l'étape de décodage. Considérez deux architectures :
\[\text{Architecture A : }\tau_i=\tau,\qquad\text{Architecture B : }\tau_i=f(s_{i-1})>0,\]
où l'architecture A utilise une unique température apprise, tandis que l'architecture B prédit la température à partir de l'état caché précédent du décodeur.
\begin{enumerate}[label=(\alph*),leftmargin=16pt]
\item Donnez une limite de l'architecture A et un avantage potentiel de l'architecture B.
\item Un collègue affirme qu'une petite valeur prédite de $\tau_i$ signifie que le modèle se concentre sur un seul état de l'encodeur, tandis qu'une grande valeur de $\tau_i$ signifie qu'il combine plusieurs états de l'encodeur. À partir de vos résultats précédents, expliquez pourquoi cette interprétation peut être incorrecte.
\end{enumerate}
\end{enumerate}
%\newpage

\Answer{

}
% Q4

\question{4 (4-2-4)}{Optimisation et régularisation}
Dans cette question, nous considérons deux approches de régularisation au cours de l'optimisation. Soit $\theta$ l'ensemble des paramètres entraînables d'un réseau de neurones, dont la fonction de perte est
\[\mathcal L(f(x^{(i)},\theta),y^{(i)}).\]
Une première approche est la \textit{décroissance des poids} (\textit{weight decay}), qui modifie la règle de mise à jour de l'optimiseur. À chaque itération, les paramètres sont mis à jour selon
\[\theta_{t+1}=\theta_t-\Delta\theta_t-\lambda\theta_t,\]
où $\Delta\theta_t$ dépend de l'optimiseur. Pour SGD, il s'agit du produit du taux d'apprentissage par le gradient. Pour RMSprop et Adam, la mise à jour dépend de moyennes mobiles. Reportez-vous aux diapositives du cours pour les règles de mise à jour correspondantes.

Une autre approche est la \textit{régularisation $L_2$}, qui modifie la fonction de perte :
\[\mathcal L_{\mathrm{reg}}=\mathcal L(f(x^{(i)},\theta),y^{(i)})+\gamma\|\theta\|_2^2.\]
\begin{enumerate}
\item Considérez l'affirmation suivante :
\begin{quote}\og Utiliser une régularisation $L_2$ avec la mise à jour habituelle de l'optimiseur est équivalent à ajouter une décroissance des poids à la mise à jour de l'optimiseur.\fg\end{quote}
Cette affirmation est-elle valable pour SGD, RMSprop et Adam ? Fournissez des arguments mathématiques pour chacune de vos réponses.
\item Comparons maintenant les régularisations $L_1$ et $L_2$. La fonction de perte régularisée par $L_1$ est
\[\mathcal L_{\mathrm{reg}}=\mathcal L(f(x^{(i)},\theta),y^{(i)})+\gamma\|\theta\|_1.\]
Considérez un modèle de régression linéaire avec régularisation $L_1$ :
\[
\hat\beta=\arg\min_{\beta}\left\{\frac12\sum_{i=1}^{N}\left(y_i-\beta_0-\sum_{j=1}^{p}x_{ij}\beta_j\right)^2+\lambda\sum_{j=1}^{p}|\beta_j|\right\}.
\]
Le modèle correspondant avec régularisation $L_2$ est
\[
\hat\beta=\arg\min_{\beta}\left\{\frac12\sum_{i=1}^{N}\left(y_i-\beta_0-\sum_{j=1}^{p}x_{ij}\beta_j\right)^2+\lambda\sum_{j=1}^{p}|\beta_j|^2\right\}.
\]
Pour deux paramètres $\beta_0$ et $\beta_1$, démontrez que les ensembles de niveau de la perte quadratique non régularisée sont des ellipses.
\item Laquelle des régularisations $L_1$ et $L_2$ favorise la parcimonie ? Donnez une explication intuitive, sans entrer dans les détails mathématiques. Quel terme de régularisation choisiriez-vous pour sélectionner un sous-ensemble de variables explicatives ?
\end{enumerate}
%\newpage

\Answer{

}
% Q5

\question{5 (5-5-5-5)}{Analyse d'article : Vision Transformer (ViT)}
Dans cette question, vous devez rédiger une analyse d'une page de l'article \url{https://arxiv.org/abs/2010.11929}, \textit{An Image Is Worth 16x16 Words: Transformers for Image Recognition at Scale} (« Une image vaut 16 × 16 mots : les Transformers pour la reconnaissance d'images à grande échelle »).

Votre analyse doit comporter les quatre sections suivantes : \textbf{Résumé}, \textbf{Points forts}, \textbf{Points faibles} et \textbf{Réflexions}. Pour chacune, nous proposons des questions directrices à prendre en compte lors de la lecture de l'article. Servez-vous-en pour structurer votre analyse, plutôt que d'y répondre une par une.

\textbf{(5.1) Résumé :}
\begin{enumerate}[label=(\alph*),leftmargin=25pt,itemsep=2pt,topsep=2pt]
\item Quel problème l'article cherche-t-il à résoudre ?
\item Quelles sont ses principales contributions ?
\item Comment ViT traite-t-il les images au moyen d'un Transformer, et en quoi cela diffère-t-il des réseaux de neurones convolutifs ?
\item Quels sont les principaux résultats expérimentaux ? Privilégiez les faits plutôt que les opinions.
\end{enumerate}
\textbf{(5.2) Points forts :}
\begin{enumerate}[label=(\alph*),leftmargin=25pt,itemsep=2pt,topsep=2pt]
\item Quels sont les principaux avantages de l'approche proposée ?
\item L'article présente-t-il une avancée conceptuelle ou empirique importante ?
\item Les expériences sont-elles bien conçues et convaincantes ?
\item Qu'est-ce qui rend cette approche utile à la communauté scientifique ?
\end{enumerate}
\textbf{(5.3) Points faibles :}
\begin{enumerate}[label=(\alph*),leftmargin=25pt,itemsep=2pt,topsep=2pt]
\item Quelles sont les principales limites de ViT ?
\item Manque-t-il des méthodes de référence, des jeux de données ou des comparaisons expérimentales ?
\item Certains choix de conception sont-ils insuffisamment justifiés ?
\item Comment pourrait-on améliorer la méthode proposée ou son évaluation ?
\end{enumerate}
\textbf{(5.4) Réflexions :}
\begin{enumerate}[label=(\alph*),leftmargin=25pt,itemsep=2pt,topsep=2pt]
\item Comment cet article se rattache-t-il aux concepts étudiés en cours, notamment les réseaux convolutifs, les mécanismes d'attention et les Transformers ?
\item Quelles pistes de recherche futures pourraient prolonger ces travaux ?
\item Quelles nouvelles idées ou expériences souhaiteriez-vous explorer ?
\end{enumerate}
Cette question est subjective : différentes réponses solidement argumentées seront acceptées. Vous devez analyser l'article de manière critique et apporter votre propre point de vue et vos propres idées, au-delà de ce qui est discuté dans l'article lui-même.
\end{document}